\documentclass[10pt,aspectratio=169]{beamer}

\usetheme[progressbar=frametitle,sectionpage=none]{arara}
\usepackage{tikz}
\usepackage{pgfplots}
\pgfplotsset{compat=1.18}

\title{A living mosaic}
\subtitle{Biodiversity across Brazil}
\author{ARARA · FIELD NOTES 01}
\institute{Amazonia · Cerrado · Atlantic Forest}
\date{November 6, 2026}

\begin{document}

\begin{frame}[plain]
  \begin{tikzpicture}[remember picture,overlay]
    \fill[araraPrimary] (current page.south west) rectangle (current
    page.north east);
    \fill[araraGold] ([xshift=-2.2cm,yshift=1.5cm]current page.east) --
    ([xshift=-4.8cm,yshift=0.5cm]current page.east) --
    ([xshift=-2.4cm,yshift=-0.4cm]current page.east) --
    ([xshift=-0.3cm,yshift=0.5cm]current page.east) -- cycle;
    \fill[araraBlue] ([xshift=-2.2cm,yshift=1.2cm]current page.east)
    circle (0.45cm);
    \fill[araraInk] ([xshift=-2.2cm,yshift=1.2cm]current page.east)
    circle (0.2cm);
    \draw[araraPaper!40,line width=0.7pt]
    ([xshift=-5.4cm,yshift=-1.6cm]current page.east) .. controls
    ([xshift=-3.7cm,yshift=-0.7cm]current page.east) and
    ([xshift=-2.8cm,yshift=-2.2cm]current page.east) ..
    ([xshift=-0.5cm,yshift=-1.1cm]current page.east);
    \fill[araraGold] ([xshift=-1.1cm,yshift=-2.2cm]current page.east)
    circle (0.08cm);
    \fill[araraPaper] ([xshift=-4.2cm,yshift=2.1cm]current page.east)
    circle (0.06cm);
  \end{tikzpicture}
  \vspace{0.6cm}
  {\color{araraGold}\small\bfseries ARARA  /  FIELD NOTES 01}\par
  \vspace{0.7cm}
  {\usebeamerfont{title}\color{araraPaper}A living\\mosaic}\par
  \vspace{0.35cm}
  {\large\color{araraPaper!80}Biodiversity across Brazil}\par
  \vfill
  {\small\color{araraPaper!65}AMAZONIA  ·  CERRADO  ·  ATLANTIC FOREST}\par
  \vspace{0.7cm}
\end{frame}

\begin{frame}{One country. Many worlds.}
  \begin{columns}[T,onlytextwidth]
    \column{0.53\textwidth}
    {\Large Brazil’s living landscapes are connected by\alert{ water,
    climate, and movement}.}\par
    \vspace{0.6cm}
    Each biome holds its own rhythm. Together, they form a network that
    sustains people and wildlife far beyond its borders.
    \column{0.42\textwidth}
    \begin{tikzpicture}[scale=0.92]
      \fill[araraGold!65] (0,0) -- (1.6,0.3) -- (2.3,1.4) --
      (1.6,2.7) -- (0.1,2.3) -- cycle;
      \draw[araraPrimary,line width=1.3pt] (0,0) -- (1.6,0.3) --
      (2.3,1.4) -- (1.6,2.7) -- (0.1,2.3) -- cycle;
      \fill[araraPrimary] (0.9,1.8) circle (0.11);
      \fill[araraBlue] (1.5,1.25) circle (0.11);
      \fill[araraGold!80!black] (0.9,0.65) circle (0.11);
      \node[anchor=west] at (2.55,1.8) {Amazonia};
      \node[anchor=west] at (2.55,1.25) {Cerrado};
      \node[anchor=west] at (2.55,0.7) {Atlantic Forest};
    \end{tikzpicture}
  \end{columns}
\end{frame}

\section{The signal}

\begin{frame}{A remarkable share of life}
  \begin{columns}[T,onlytextwidth]
    \column{0.48\textwidth}
    \begin{block}{A global hotspot}
      Brazil is home to roughly one in ten of the world’s known species.
    \end{block}
    \vspace{0.4cm}
    \small Species counts keep changing as researchers document life in
    remote forests, wetlands, and grasslands.
    \column{0.48\textwidth}
    \centering
    \begin{tikzpicture}
      \fill[araraGold!28] (0,0) circle (1.35);
      \fill[araraPrimary] (0,0) circle (1.35);
      \fill[araraGold] (0,0) arc (90:54:1.35) -- (0,0) -- cycle;
      \node[align=center,text=white] at (0,-0.05)
      {{\fontsize{34}{36}\selectfont\bfseries 10\%}\\[-1mm]\small of
      known species};
    \end{tikzpicture}
  \end{columns}
\end{frame}

\begin{frame}{What helps species thrive?}
  \begin{itemize}
    \item Protect large, connected habitats.
    \item Restore native plants along rivers and forest edges.
    \item Work with Indigenous and local communities.
  \end{itemize}
\end{frame}

\begin{frame}{A connected landscape}
  \framesubtitle{Habitat corridors link protected areas}
  \begin{figure}
    \centering
    \begin{tikzpicture}[scale=0.9]
      \fill[araraBlue!15] (-5,-1.1) rectangle (5,1.1);
      \fill[araraPrimary] (-4.6,-0.8) rectangle (-2.6,0.8);
      \fill[araraPrimary] (2.6,-0.8) rectangle (4.6,0.8);
      \draw[araraGold,line width=8pt] (-2.6,0) .. controls (-1.4,0.9)
      and (1.4,-0.9) .. (2.6,0);
      \node[text=white,font=\small\bfseries] at (-3.6,0) {Reserve A};
      \node[text=white,font=\small\bfseries] at (3.6,0) {Reserve B};
      \node[araraInk,font=\small\bfseries] at (0,1.55) {RIVER CORRIDOR};
    \end{tikzpicture}
    \caption{A corridor lets wildlife move between protected habitats.}
  \end{figure}
\end{frame}

\begin{frame}{Local action, regional impact}
  \begin{columns}[T,onlytextwidth]
    \column{0.47\textwidth}
    \begin{block}{At the site}
      Restore native vegetation, protect riverbanks, and reduce fire risk.
    \end{block}
    \column{0.47\textwidth}
    \begin{exampleblock}{Across the region}
      Connected habitats support wildlife movement and healthier water flows.
    \end{exampleblock}
  \end{columns}
\end{frame}

\begin{frame}{A simple water balance}
  \framesubtitle{Rainfall, evaporation, and river flow shape wetlands}
  \begin{equation*}
    \Delta S = P - E - Q
  \end{equation*}
  \begin{align*}
    P & = \text{precipitation} \\
    E & = \text{evaporation} \\
    Q & = \text{water flowing out}
  \end{align*}
  \begin{center}\small $\Delta S$ is the change in stored water.
  \end{center}
\end{frame}

\begin{frame}{Life moves with the seasons}
  \begin{columns}[T,onlytextwidth]
    \column{0.42\textwidth}
    Flood and drought reshape the Pantanal through the year.
    \vspace{0.5cm}
    \begin{exampleblock}{A moving habitat}
      Rising water opens feeding grounds. Receding water gathers wildlife
      around shrinking wetlands.
    \end{exampleblock}
    \column{0.55\textwidth}
    \begin{tikzpicture}
      \begin{axis}[
          width=\linewidth,height=5cm,
          axis lines=left,axis line style={draw=araraPrimary!45},
          tick style={draw=none},
          xtick={1,2,3,4},xticklabels={Jan,Apr,Jul,Oct},
          ytick=\empty,ymin=0,ymax=10,
          xlabel={Season},ylabel={Flood pulse},
          xlabel style={font=\small},ylabel style={font=\small},
          xticklabel style={font=\small,text=araraInk!65},
          clip=false
        ]
        \addplot[araraBlue,line width=2.5pt,smooth,mark=*,mark size=2.5pt]
        coordinates {(1,2)(2,7)(3,9)(4,3)};
        \addplot[araraBlue!18,draw=none,smooth]
        coordinates {(1,0)(1,2)(2,7)(3,9)(4,3)(4,0)} -- cycle;
      \end{axis}
    \end{tikzpicture}
  \end{columns}
\end{frame}

\section{The pressure}

\begin{frame}{A changing edge}
  \begin{columns}[T,onlytextwidth]
    \column{0.58\textwidth}
    \begin{tikzpicture}
      \begin{axis}[
          width=\linewidth,height=5.2cm,
          ybar,bar width=18pt,
          axis lines=left,axis line style={draw=araraPrimary!45},
          tick style={draw=none},
          symbolic x coords={Forest,Crops,Pasture},xtick=data,
          ytick=\empty,ymin=0,ymax=100,
          nodes near coords, nodes near coords
          style={font=\small,text=araraInk},
          xticklabel style={font=\small,text=araraInk!70},
          enlarge x limits=0.3
        ]
        \addplot[fill=araraPrimary,draw=none] coordinates
        {(Forest,84)(Crops,61)(Pasture,48)};
      \end{axis}
    \end{tikzpicture}
    \column{0.36\textwidth}
    \textbf{Change travels.}\\[0.2cm]
    Habitat loss can alter fire patterns, water flows, and species movement
    across an entire landscape.
    \vspace{0.45cm}
    {\scriptsize Illustrative landscape index, not measured field data.}
  \end{columns}
\end{frame}

\begin{frame}{Protection works at landscape scale}
  \begin{columns}[T,onlytextwidth]
    \column{0.31\textwidth}
    \begin{block}{Connect}
      Keep habitat corridors open between protected areas.
    \end{block}
    \column{0.31\textwidth}
    \begin{exampleblock}{Restore}
      Bring native vegetation back along rivers and degraded land.
    \end{exampleblock}
    \column{0.31\textwidth}
    \begin{alertblock}{Partner}
      Back Indigenous stewardship and local conservation work.
    \end{alertblock}
  \end{columns}
  \vfill
  \begin{center}\Large\bfseries A living mosaic needs every piece.
  \end{center}
\end{frame}

\begin{frame}{Read change at two scales}
  \framesubtitle{Two views of the same living system}
  \begin{columns}[T,onlytextwidth]
    \column{0.47\textwidth}
    \textbf{On the ground}\par
    \vspace{0.25cm}
    Local choices shape habitat quality: how land is managed, where native
    plants return, and whether riverside forests stay connected.
    \column{0.47\textwidth}
    \textbf{Across the region}\par
    \vspace{0.25cm}
    Many local changes add up. Together, they influence wildlife movement,
    water cycles, and resilience across whole biomes.
  \end{columns}
\end{frame}

\begin{frame}[plain]
  \begin{tikzpicture}[remember picture,overlay]
    \fill[araraInk] (current page.south west) rectangle (current
    page.north east);
    \fill[araraGold] ([xshift=-1.8cm,yshift=1cm]current page.east) --
    ([xshift=-3cm,yshift=0.45cm]current page.east) --
    ([xshift=-1.8cm,yshift=-0.1cm]current page.east) -- cycle;
    \fill[araraBlue] ([xshift=-2.35cm,yshift=0.45cm]current
    page.east) circle (0.13cm);
  \end{tikzpicture}
  \vspace{1.3cm}
  {\color{araraGold}\small\bfseries ARARA  /  FIELD NOTES 01}\par
  \vspace{0.6cm}
  {\usebeamerfont{title}\color{araraPaper}Keep it wild.}\par
  \vspace{0.4cm}
  {\large\color{araraPaper!75}Protect the links that let life thrive.}\par
\end{frame}

\end{document}
